\DocumentMetadata{
  pdfversion=2.0,
  pdfstandard=ua-2,
  lang=en-US,
  tagging=on,
  tagging-setup={math/setup=mathml-SE,tabsorder=structure}
}
\documentclass[11pt,letterpaper]{article}
\usepackage[margin=0.68in,includefoot]{geometry}
\usepackage{newcomputermodern}
\usepackage{amsmath,amscd,mathtools}
\usepackage{tikz,adjustbox}
\providecommand{\square}{\mdlgwhtsquare}
\usepackage[hidelinks]{hyperref}
\hypersetup{
  pdftitle={Algebra PhD Exam, September 1990},
  pdfauthor={University of Florida Department of Mathematics},
  pdfsubject={Graduate mathematics examination},
  pdfdisplaydoctitle=true,
  bookmarksopen=true
}
\setcounter{secnumdepth}{0}
\setlength{\parindent}{0pt}
\setlength{\parskip}{0.28em}
\setlength{\emergencystretch}{3em}
\setlength{\leftmargini}{2.15em}
\setlength{\leftmarginii}{2.65em}
\setlength{\leftmarginiii}{3em}
\renewcommand{\labelenumi}{\textbf{\theenumi.}}
\pagestyle{empty}
\raggedbottom
\allowdisplaybreaks
\newenvironment{examproblems}[1]{%
  \begin{enumerate}%
  \setcounter{enumi}{#1}%
  \setlength{\topsep}{0.35em}%
  \setlength{\partopsep}{0pt}%
  \setlength{\itemsep}{0.48em}%
  \setlength{\parsep}{0.18em}%
}{\end{enumerate}}
\newenvironment{examsubparts}{%
  \begin{enumerate}%
  \setlength{\topsep}{0.18em}%
  \setlength{\partopsep}{0pt}%
  \setlength{\itemsep}{0.16em}%
  \setlength{\parsep}{0.1em}%
}{\end{enumerate}}
\newenvironment{examinstructions}{%
  \par\begingroup\itshape
}{%
  \par\endgroup\vspace{0.18em}
}
\makeatletter
\renewcommand\section{\@startsection{section}{1}{0pt}%
  {-1.25ex plus -.25ex minus -.1ex}%
  {0.55ex plus .1ex}%
  {\normalfont\large\bfseries}}
\renewcommand{\@maketitle}{%
  \newpage\null\vskip -1em%
  {\footnotesize\bfseries
    \href{https://www.ufl.edu/}{University of Florida}, \href{https://math.ufl.edu/}{Department of Mathematics}\par}%
  \vskip -0.5em%
  \tagstructbegin{tag=Title}%
    {\LARGE\bfseries\href{https://gma.math.ufl.edu/exams/algebra-phd/}{Algebra PhD Exam}, September 1990\par}%
  \tagstructend%
  \vskip 0.25em%
  \tagmcbegin{artifact}%
    \hrule height 1pt%
  \tagmcend%
  \vskip 1em%
}
\makeatother
\title{Algebra PhD Exam, September 1990}
\author{}
\date{}
\begin{document}
\maketitle
\thispagestyle{empty}
\begin{examinstructions}
Read through the questions carefully before you begin. Your six best solutions will count. Theorems from the course may be quoted (within reason) as long as you give a clear indication and state the result in full.
\end{examinstructions}
\begin{examproblems}{0}
\item
Let~\(P\) be a nontrivial, finite~\(p\)-group, and~\(Q\) a subgroup. Consider the action of~\(P\) on the set \(P/Q\) of left cosets of~\(Q\) by left multiplication: 
\[
P\times P/Q\longrightarrow P/Q,\qquad (p,xQ)\longmapsto (px)Q.
\]
 \(Q\) acts on \(P/Q\) by restriction.
\begin{examsubparts}
\item[\textnormal{a)}]
Show that \(xQ\) is a fixed point of~\(Q\) if and only if \(x\in N_P(Q)=\{p\in P\mid pQp^{-1}=Q\}\).
\item[\textnormal{b)}]
Deduce that if~\(Q\) is a proper subgroup of~\(P\) then~\(Q\) is also a proper subgroup of \(N_P(Q)\).
\item[\textnormal{c)}]
Hence show that every maximal subgroup of~\(P\) is normal of index~\(p\).
\end{examsubparts}
\item
Let~\(A\) be a Dedekind domain with field of fractions~\(K\). Show that every ideal of~\(A\) can be generated by two elements. (Hint: use the Chinese remainder theorem.)
\item
Prove that any simple group of order 60 is isomorphic to \(A_5\). (you need not show that \(A_5\) is simple).
\item
Let~\(A\), \(B\) and~\(C\) be rings (with 1) and suppose we have a commutative diagram of ring homomorphisms: 
\[
\begin{array}{ccc}
A & \xrightarrow{\alpha} & B \\
& {}_{\gamma}\searrow \quad \swarrow_{\beta} & \\
& C &
\end{array}
\]
\begin{examsubparts}
\item[\textnormal{a)}]
Define the structure of an \((A,B)\) bimodule on~\(B\), ie make~\(B\) into a left~\(A\)-module and a right~\(B\)-module such that \((a\cdot x)b=a(x\cdot b)\) for \(a\in A\), \(x,b\in B\).
\item[\textnormal{b)}]
Show that for any right~\(A\)-module~\(M\), we have a natural isomorphism 
\[
M\otimes_A C\cong (M\otimes_A B)\otimes_B C.
\]
 (Remember what ‘natural’ means in terms of the functors \(-\otimes_A C\) and \((-\otimes_A B)\otimes_B C\).)
\end{examsubparts}
\item
Let \(\{\sigma_i\}_{i=1}^n\) be a set of~\(n\) distinct automorphisms of a field~\(k\). Show that if for \(a_1,\ldots,a_n\in k\), 
\[
\sum a_i\sigma_i(u)=0\qquad\text{for all }u\in k,
\]
 then \(a_i=0\) for all~\(i\). (Hint: show that there is no ‘shortest’ relation with \(a_i\ne0\).)
\item
Give examples of the following (with proof).
\begin{examsubparts}
\item[\textnormal{a)}]
An irreducible cubic polynomial with rational coefficients whose Galois group has order 3.
\item[\textnormal{b)}]
Two normal extensions \(K\subset L\), \(L\subset M\) such that extension \(K\subset M\) is not normal.
\end{examsubparts}
\item
Let~\(R\) be a ring and~\(M\) simple (irreducible)~\(R\)-module.
\begin{examsubparts}
\item[\textnormal{a)}]
Prove Schur's Lemma: every nonzero~\(R\)-module map from~\(M\) to itself is an isomorphism.
\item[\textnormal{b)}]
Give an example to show that the converse of Schur's Lemma is false.
\end{examsubparts}
\item
Let \(\mathbb Z_p\) denote the~\(p\)-adic integers and for \(n\in\mathbb N\), 
\[
\varepsilon_n:\mathbb Z_p\longrightarrow \mathbb Z/p^n\mathbb Z,
\]
 the map defining the structure of \(\mathbb Z_p\) as \(\varprojlim_n\mathbb Z/p^n\mathbb Z\). Let \(U_n=1+p^n\mathbb Z_p\).
\begin{examsubparts}
\item[\textnormal{a)}]
Show that \(U_n\) is a subgroup of the group~\(U\) of units of \(\mathbb Z_p\).
\item[\textnormal{b)}]
Show that \(U_n\) is the kernel of \(\varepsilon_n|_U:U\to(\mathbb Z/p^n\mathbb Z)^*\) and that \(U/U_1\cong\mathbb F_p^*\) is cyclic of order \(p-1\).
\item[\textnormal{c)}]
For \(n\ge1\) show that the map 
\[
1+p^nx\longmapsto x\text{ modulo }p
\]
 defines an isomorphism 
\[
U_n/U_{n+1}\xrightarrow{\sim}\mathbb Z/p\mathbb Z.
\]
\item[\textnormal{d)}]
Deduce that \(U_1/U_n\) is a~\(p\)-group and that \(U/U_n\) has a unique subgroup of order \(p-1\).
\item[\textnormal{e)}]
Prove that~\(U\) has a subgroup of order \(p-1\). (Hint: \(U\cong\varprojlim U/U_n\).)
\end{examsubparts}
\item
Let~\(D\) be a division ring and~\(M\) a~\(D\)-module. A subring~\(R\) of \(\operatorname{End}_D(M)\) is said to be dense if for any natural number~\(n\) and \(v_1,\ldots,v_n\in M\) which are linearly independent over~\(D\), and any~\(n\) elements \(w_1,\ldots,w_n\) of~\(M\), there exists \(r\in R\) such that \(w_i=r(v_i)\), \(i=1,\ldots,n\).
\begin{examsubparts}
\item[\textnormal{a)}]
Show that if \(\dim_D M\) is finite then a dense subring of \(\operatorname{End}_D M\) must be equal to \(\operatorname{End}_D M\).
\item[\textnormal{b)}]
Show that if~\(M\) is infinite-dimensional over~\(D\) and~\(R\) is a dense subring of \(\operatorname{End}_D M\), then for every natural number~\(m\), \(R\) has a subring \(S_m\) which maps homomorphically onto \(\operatorname{Mat}_m(D)\).
\end{examsubparts}
\item
This problem has three parts.
\begin{examsubparts}
\item[\textnormal{a)}]
State Hilbert's Nullstellensatz, giving the main definitions involved.
\item[\textnormal{b)}]
Let~\(k\) be an algebraically closed field. Show that every algebraic subset of \(k^n\) can be decomposed uniquely as a finite union of (irreducible) varieties.
\item[\textnormal{c)}]
Give an example to show that in b) the union may not be disjoint.
\end{examsubparts}
\end{examproblems}
\end{document}
